% ============================================================================
%  CHƯƠNG 4b. THÍ NGHIỆM, ĐÁNH GIÁ VÀ KẾT QUẢ
% ============================================================================
\clearpage
\section{Thí nghiệm, đánh giá và kết quả}
\label{sec:results}

\begin{warnbox}
\textbf{Trạng thái.} Báo cáo đang ở giai đoạn đề cương: thiết kế thí nghiệm và cách đánh giá đã chốt, mã nguồn
(gói \code{nodulecxr}) đã cài đặt xong và qua 75 kiểm thử trên CPU; các bước cần GPU chưa chạy, nên chưa có kết quả. Các ô \textcolor{badred}{[--]} sẽ được điền tự
động từ file dự đoán sau khi chạy E0--E5.
\end{warnbox}

\subsection{Thiết kế thí nghiệm}

Bảng~\ref{tab:experiments} liệt kê sáu thí nghiệm; Hình~\ref{fig:workflow} cho thứ tự chạy và các phụ thuộc.
Các thí nghiệm phải chạy đúng thứ tự này. E4 và E5 phụ thuộc vào cấu hình được chọn, và cấu hình đó chỉ được chọn
từ \code{val\_selection.csv}, file do \code{make\_report.py} sinh ra sau E3. Lệnh cụ thể của từng bước nằm trong
\code{code/RUNBOOK.md}; Phụ lục~\ref{app:protocol} tóm tắt đầu vào, đầu ra và bảng mà mỗi bước điền.

\begin{table}[H]
\centering
\caption{Các thí nghiệm.}
\label{tab:experiments}
\small
\begin{tabularx}{\textwidth}{@{}l L L@{}}
\toprule
\textbf{Mã} & \textbf{Nội dung} & \textbf{Kết quả cần có} \\
\midrule
E0 & Chuẩn bị dữ liệu: số ca Nodule/Mass, phân bố kích thước box, độ đồng thuận giữa bác sĩ; hợp nhất box & Bảng~\ref{tab:data-stats} \\
E1 & Baseline B1, B2 trên VinBigData, đánh giá trên 3 tập test, mỗi loại 3 seed & AUC, sens@spec0,90, FROC \\
E2 & Zero-shot C0 cho MedGemma-4B và Qwen2.5-VL-3B: 1 lần nhiệt độ 0 (dòng chính) và 2 lần nhiệt độ 0,7 (đo độ ổn định) & Bảng zero-shot, tỷ lệ parse lỗi, tỷ lệ đồng thuận \\
E3 & LoRA C1, C2, C3 trên MedGemma-4B, mỗi cấu hình 3 seed & Bảng so sánh, đường học, sensitivity theo epoch \\
E4 & C4: cấu hình tốt nhất áp lên Qwen2.5-VL-3B, 1--3 seed & So sánh mô hình y khoa với mô hình đa năng \\
E5 & Kiểm toán ``nhìn đúng chỗ'' (Mục~\ref{sec:audit}) cho C0, cấu hình LoRA tốt nhất và B2; phân tầng & Bảng đổi đáp án, sensitivity theo tầng \\
\bottomrule
\end{tabularx}
\end{table}

\begin{figure}[H]
\centering
\begin{tikzpicture}[x=1cm,y=1cm, fb/.append style={font=\scriptsize, minimum height=0.85cm}]
  \node[fgray, text width=1.5cm, minimum height=1.15cm] (pc) at (0.9,0) {\textbf{Kiểm tra trước}\\(7 mục)};
  \node[fb, text width=1.5cm, minimum height=1.15cm] (e0) at (3.1,0) {\textbf{E0}\\dữ liệu, manifest};
  \node[fgreen, text width=1.8cm] (e1) at (5.6,1.35) {\textbf{E1} B1, B2\\3 seed};
  \node[fgray, text width=1.8cm] (e2) at (5.6,0) {\textbf{E2} C0\\run 0, 1, 2};
  \node[fblue, text width=1.8cm] (e3) at (5.6,-1.35) {\textbf{E3} C1--C3\\3 seed};
  \node[forange, text width=1.9cm, minimum height=1.15cm] (mr1) at (8.35,0) {\code{make\_report}\\(lần 1)};
  \node[forange, text width=2.0cm, minimum height=1.15cm] (sel) at (10.9,0) {\code{val\_selection}\\chọn theo \code{val\_sens90}};
  \node[fblue, text width=2.0cm] (e4) at (13.85,1.0) {\textbf{E4} C4\\Qwen + LoRA};
  \node[fred, text width=2.0cm] (e5) at (13.85,-1.0) {\textbf{E5} kiểm toán\\T1, T2, T3};
  \node[fb, text width=4.8cm] (mr2) at (10.4,-2.75) {\code{make\_report} (lần cuối) $\rightarrow$ \code{summary.md}\\và mọi bảng của báo cáo};
  \draw[farr] (pc) -- (e0);
  \draw[farr] (e0.east) -- ++(0.3,0) |- (e1.west);
  \draw[farr] (e0) -- (e2);
  \draw[farr] (e0.east) -- ++(0.3,0) |- (e3.west);
  \draw[farr] (e1.east) -- ++(0.3,0) |- (mr1.west);
  \draw[farr] (e2) -- (mr1);
  \draw[farr] (e3.east) -- ++(0.3,0) |- (mr1.west);
  \draw[farr] (mr1) -- (sel);
  \draw[farr] (sel.east) -- ++(0.3,0) |- (e4.west);
  \draw[farr] (sel.east) -- ++(0.3,0) |- (e5.west);
  \draw[farr] (e5.south) |- (mr2.east);
  \draw[farr] (e4.east) -- ++(0.25,0) |- ([yshift=0.15cm]mr2.east);
  \draw[farr] (e3.south) |- (mr2.west);
  \begin{scope}[on background layer]
    \node[draw=warnorange, dashed, thick, rounded corners=4pt, fill=warnorange!4, inner sep=6pt,
          fit=(mr1)(sel), label={[font=\scriptsize\bfseries, text=warnorange]above:CỔNG: chỉ đọc cột validation}] {};
  \end{scope}
  \node[fnote, text=badred, align=left, anchor=north west, text width=4.0cm] at (0.0,-1.0)
     {Tập test được dự đoán cùng lúc nhưng \textbf{không} dùng để chọn cấu hình hay ngưỡng; số trên test chỉ đọc
     một lần ở bảng cuối.};
\end{tikzpicture}
\caption{Thứ tự chạy thí nghiệm. E1--E3 độc lập với nhau sau E0. \code{make\_report.py} lần 1 sinh
\code{val\_selection.csv} và \code{thresholds.json}; cổng chọn (khung cam) chỉ dùng số trên validation để chọn cấu
hình LoRA tốt nhất cho E4 và E5. Lần chạy cuối của \code{make\_report.py} gộp mọi file dự đoán thành bảng báo cáo.}
\label{fig:workflow}
\end{figure}

\subsection{Chỉ số đánh giá}

\paragraph{Phát hiện mức ảnh.}
\begin{itemize}
  \item Sensitivity, specificity, PPV, NPV, F1; AUC-ROC và AUC-PR.
  \item \textbf{Chỉ số chính:} sensitivity ở specificity 0,90 (và 0,95). Ngưỡng $t_{90}$, $t_{95}$ chọn trên
        validation, riêng cho từng bộ (cấu hình, seed, lần chạy), rồi áp nguyên lên cả ba tập test.
        \textbf{Tiêu chí thành công:} sensitivity $\ge 0{,}60$ ở specificity 0,90 trên VinBigData test.
  \item \textbf{Chỉ số trên văn bản sinh ra} (chỉ có với VLM chạy \code{--generate all}):
        \code{raw\_sens}/\code{raw\_spec} tính trên mọi ảnh, câu trả lời không parse được tính là ``no''. Các chỉ số
        chế độ báo bình thường, \code{no\_on\_pos} (tỷ lệ ``no'' trên ca dương) và \code{yes\_on\_neg} (tỷ lệ ``yes''
        trên ca âm), chỉ tính trên câu trả lời parse được. Tỷ lệ không parse được (\code{unparsed}) báo riêng.
        Mọi chỉ số này được tính lại từ \code{raw\_text} bằng parser hiện hành.
  \item \textbf{C0:} dòng chính chỉ dùng run 0 (nhiệt độ 0, tối đa 256 token). Run 1--2 (nhiệt độ 0,7) chỉ dùng cho
        tỷ lệ đồng thuận trong \code{run\_agreement.csv}, tính trên các ảnh parse được ở mọi lần chạy.
  \item \textbf{Theo tỷ lệ bệnh}~\cite{mosquera2024imbalance,kim2023multicentre}: ECE, Brier và balanced Brier;
        PPV quy về tỷ lệ bệnh 1\% và 3\%. Trên JSRT tỷ lệ nốt là 62\%, nên PPV ở đó không có ý nghĩa lâm sàng.
\end{itemize}

\paragraph{Định vị.} IoU trung bình và tỷ lệ ``đúng và IoU $\ge 0{,}3$'' (ngưỡng thấp vì nốt nhỏ; báo thêm ở 0,5)
chỉ dùng các box có điểm $\ge t_{90}$, ngưỡng ứng với specificity 0,90 (box của VLM mang điểm của ảnh). FROC dùng
mọi box, áp cho B2 và các VLM có box. Trên JSRT, chỉ xét box có điểm cao nhất; ``trúng'' khi tâm box nằm trong
bán kính nốt cộng biên 5\,mm (12,5\,px trên ảnh 896); báo thêm khoảng cách tâm trung vị (mm, quy đổi 0,4\,mm/px),
theo cách của NoduLoCC~\cite{mustafic2026noduloc}.

\paragraph{Phân tầng.} Sensitivity được báo theo tầng kích thước (\code{strata.csv}): VinBigData và NIH theo
tertile diện tích box tương đối, ngưỡng tertile lấy từ ca dương train VinBigData (\code{e0\_data\_edges.json});
JSRT theo mm: $\le 10$, 10--20, $> 20$\,mm; thêm theo độ mờ JSRT 1--5.

\paragraph{Thống kê.} Trung bình $\pm$ độ lệch chuẩn qua các seed; CI 95\% bằng bootstrap theo ảnh (1.000 lần lấy
mẫu lại), phân tầng theo nhãn (CI theo từng seed lưu trong \code{summary\_per\_seed.csv}); McNemar cho sens/spec,
DeLong cho AUC, hiệu chỉnh Holm khi so nhiều cặp (\code{comparisons.csv}). Bootstrap lấy mẫu theo ảnh, không gộp
theo bệnh nhân, nên CI có thể hẹp hơn thực tế nếu một bệnh nhân có nhiều ảnh; đây là một hạn chế.
Hình~\ref{fig:evalflow} tóm tắt luồng đánh giá.

\begin{figure}[H]
\centering
\begin{tikzpicture}[x=1cm,y=1cm, fb/.append style={font=\scriptsize, minimum height=0.85cm}]
  \node[fb, text width=2.1cm, minimum height=1.3cm] (pr) at (1.15,0) {File dự đoán\\\code{preds/*.jsonl}\\+ manifest};
  \node[forange, text width=2.2cm, minimum height=1.3cm] (th) at (4.1,0) {\textbf{Validation}\\$t_{90}$, $t_{95}$ theo\\(config, seed, run)};
  \node[fb, text width=3.0cm, minimum height=1.3cm] (mt) at (7.75,0) {\textbf{Ba tập test}\\sens, spec, AUC, AUC-PR,\\ECE, Brier, PPV ở 1\% và 3\%};
  \node[fb, text width=3.0cm] (loc) at (7.75,-1.75) {\textbf{Định vị}: box có $s \ge t_{90}$\\FROC: mọi box};
  \node[fblue, text width=2.5cm] (bs) at (11.55,1.25) {Bootstrap theo ảnh\\CI 95\% (1.000 lần)};
  \node[fblue, text width=2.5cm] (mc) at (11.55,0) {McNemar (sens/spec)\\DeLong (AUC)};
  \node[fblue, text width=2.5cm] (ho) at (11.55,-1.25) {Holm khi so\\nhiều cặp};
  \node[fgreen, text width=1.4cm, minimum height=1.3cm] (out) at (14.5,0) {\code{summary}\\\code{.csv/.md}};
  \draw[farr] (pr) -- (th); \draw[farr] (th) -- node[fnote, above] {áp nguyên} (mt);
  \draw[farr] (th.south) |- (loc.west);
  \draw[farr] (mt.east) -- ++(0.35,0) |- (bs.west);
  \draw[farr] (mt) -- (mc);
  \draw[farr] (mc) -- (ho);
  \draw[farr] (loc.east) -| ([xshift=0.25cm]out.south);
  \draw[farr] (bs.east) -- ++(0.3,0) |- ([yshift=0.2cm]out.west);
  \draw[farr] (ho.east) -- ++(0.3,0) |- ([yshift=-0.2cm]out.west);
\end{tikzpicture}
\caption{Luồng đánh giá và thống kê trong \code{make\_report.py}. Ngưỡng chỉ chọn trên validation rồi áp nguyên lên
test; chỉ số định vị dùng box có điểm $\ge t_{90}$; khoảng tin cậy bằng bootstrap theo ảnh; so sánh bắt cặp
bằng McNemar và DeLong, hiệu chỉnh Holm.}
\label{fig:evalflow}
\end{figure}

\subsection{Kết quả chính}

\begin{table}[H]
\centering
\caption[Sensitivity ở specificity 0,90 và AUC-ROC trên ba tập test]{Sensitivity ở specificity 0,90 và AUC-ROC trên ba tập test, trung bình $\pm$ độ lệch chuẩn qua seed
(điền từ \code{outputs/report/summary.md} sau E1--E4; C0 chỉ dùng run 0). Mục tiêu: sens@spec0,90 $\ge 0{,}60$ trên
VinBigData test.}
\label{tab:main}
\small
\begin{tabular}{@{}l cc cc cc@{}}
\toprule
& \multicolumn{2}{c}{\textbf{VinBigData}} & \multicolumn{2}{c}{\textbf{NIH}} & \multicolumn{2}{c}{\textbf{JSRT}} \\
\cmidrule(lr){2-3}\cmidrule(lr){4-5}\cmidrule(lr){6-7}
\textbf{Cấu hình} & Sens@0,90 & AUC & Sens@0,90 & AUC & Sens@0,90 & AUC \\
\midrule
C0 MedGemma-4B   & \textcolor{badred}{[--]} & \textcolor{badred}{[--]} & \textcolor{badred}{[--]} & \textcolor{badred}{[--]} & \textcolor{badred}{[--]} & \textcolor{badred}{[--]} \\
C0 Qwen2.5-VL-3B & \textcolor{badred}{[--]} & \textcolor{badred}{[--]} & \textcolor{badred}{[--]} & \textcolor{badred}{[--]} & \textcolor{badred}{[--]} & \textcolor{badred}{[--]} \\
C1 LoRA, phân loại & \textcolor{badred}{[--]} & \textcolor{badred}{[--]} & \textcolor{badred}{[--]} & \textcolor{badred}{[--]} & \textcolor{badred}{[--]} & \textcolor{badred}{[--]} \\
C2 LoRA, + box     & \textcolor{badred}{[--]} & \textcolor{badred}{[--]} & \textcolor{badred}{[--]} & \textcolor{badred}{[--]} & \textcolor{badred}{[--]} & \textcolor{badred}{[--]} \\
C3 LoRA, + box, cân bằng & \textcolor{badred}{[--]} & \textcolor{badred}{[--]} & \textcolor{badred}{[--]} & \textcolor{badred}{[--]} & \textcolor{badred}{[--]} & \textcolor{badred}{[--]} \\
C4 Qwen2.5-VL-3B LoRA & \textcolor{badred}{[--]} & \textcolor{badred}{[--]} & \textcolor{badred}{[--]} & \textcolor{badred}{[--]} & \textcolor{badred}{[--]} & \textcolor{badred}{[--]} \\
\midrule
B1 DenseNet-121 & \textcolor{badred}{[--]} & \textcolor{badred}{[--]} & \textcolor{badred}{[--]} & \textcolor{badred}{[--]} & \textcolor{badred}{[--]} & \textcolor{badred}{[--]} \\
B2 YOLOv8       & \textcolor{badred}{[--]} & \textcolor{badred}{[--]} & \textcolor{badred}{[--]} & \textcolor{badred}{[--]} & \textcolor{badred}{[--]} & \textcolor{badred}{[--]} \\
\bottomrule
\end{tabular}
\end{table}

\paragraph{Mốc tham chiếu từ tài liệu.} Zero-shot trên JSRT, MedGemma-4b-it đạt sensitivity 5,2\% ở specificity
100\% và LLaVA-Rad đạt 27,3\% ở specificity 71,0\%~\cite{nishio2026zeroshot}. Các số này dùng đầu ra nhị phân,
không ở điểm vận hành specificity 0,90, nên chỉ để định hướng.

\subsection{Định vị}

\begin{table}[H]
\centering
\caption[Chỉ số định vị]{Chỉ số định vị (điền từ \code{summary.md} sau E1--E4). IoU và Hit chỉ tính trên box có điểm $\ge t_{90}$;
JSRT dùng box điểm cao nhất, ``trúng'' khi tâm cách tâm nốt $\le$ bán kính nốt + 5\,mm.}
\label{tab:loc}
\small
\begin{tabular}{@{}l ccc c@{}}
\toprule
& \multicolumn{3}{c}{\textbf{VinBigData test}} & \textbf{JSRT} \\
\cmidrule(lr){2-4}\cmidrule(lr){5-5}
\textbf{Cấu hình} & IoU TB & Hit@IoU0,3 & Hit@IoU0,5 & Khoảng cách tâm trung vị (mm) \\
\midrule
C0 Qwen2.5-VL-3B & \textcolor{badred}{[--]} & \textcolor{badred}{[--]} & \textcolor{badred}{[--]} & \textcolor{badred}{[--]} \\
C2 / C3 (tốt nhất) & \textcolor{badred}{[--]} & \textcolor{badred}{[--]} & \textcolor{badred}{[--]} & \textcolor{badred}{[--]} \\
C4 & \textcolor{badred}{[--]} & \textcolor{badred}{[--]} & \textcolor{badred}{[--]} & \textcolor{badred}{[--]} \\
B2 YOLOv8 & \textcolor{badred}{[--]} & \textcolor{badred}{[--]} & \textcolor{badred}{[--]} & \textcolor{badred}{[--]} \\
\bottomrule
\end{tabular}
\end{table}

\subsection{Chế độ báo bình thường và tác dụng của cân bằng lớp}

So sánh C2 với C3 theo \code{no\_on\_pos}, \code{yes\_on\_neg}, \code{unparsed}, sensitivity theo epoch (dự đoán
\code{--config <cfg>-ep<k>} trên val cho từng checkpoint của seed 0) và calibration (ECE, balanced Brier).
\textcolor{badred}{[Điền sau E3: bảng từ \code{summary.md} và đường học.]}

\subsection{Kiểm toán ``nhìn đúng chỗ''}

\begin{table}[H]
\centering
\caption[Tỷ lệ đổi đáp án dưới các phép thử che vùng]{Tỷ lệ đổi đáp án dưới các phép thử che vùng (điền từ \code{outputs/audit/*.json} sau E5; một bảng cho
\code{test}, một cho \code{test\_jsrt}).}
\label{tab:audit-res}
\small
\begin{tabular}{@{}l cccc@{}}
\toprule
\textbf{Cấu hình} & T1 (\code{flip\_t1}) & T2 (\code{flip\_t2}) & $\Delta_{\text{occ}}$ ($n_{\text{T2}}$) & T3: tỷ lệ ``yes'' \\
\midrule
C0 MedGemma-4B & \textcolor{badred}{[--]} & \textcolor{badred}{[--]} & \textcolor{badred}{[--]} & \textcolor{badred}{[--]} \\
LoRA tốt nhất  & \textcolor{badred}{[--]} & \textcolor{badred}{[--]} & \textcolor{badred}{[--]} & \textcolor{badred}{[--]} \\
B2 YOLOv8      & \textcolor{badred}{[--]} & \textcolor{badred}{[--]} & \textcolor{badred}{[--]} & \textcolor{badred}{[--]} \\
\bottomrule
\end{tabular}
\end{table}

\textcolor{badred}{[Điền sau E5 từ \code{strata.csv} và khoá \code{by\_size}/\code{by\_subtlety} của \code{audit/*.json}: biểu đồ sensitivity theo tầng kích thước (tertile diện tích với VinBigData/NIH;
$\le 10$, 10--20, $> 20$\,mm với JSRT) và theo độ mờ JSRT 1--5.]}
